\subsection{Mathematical Induction and Well-Ordering}

\begin{assumption} \label{assump:well-ordering-positive-integers} \hlt{Well-Ordering of $\N$}\\
Throughout these notes, $(\forall S\subseteq\N)\,[S\neq\emptyset\Rightarrow(\exists m\in S)(\forall n\in S)\,m\leq n]$.
\end{assumption}

\begin{theorem} \label{thm:mathematical-induction} \hlt{Mathematical Induction}\\
Let $(P_n)_{n\in\N}$ be mathematical statements. Suppose the following conditions hold.
\begin{enumerate}[label=(\roman*)]
\setlength{\itemsep}{0pt}
\item \hlt{Basis}: $P_1$ is true.
\item \hlt{Induction step}: $(\forall n\in\N)\,(P_n\Rightarrow P_{n+1})$.
\end{enumerate}
Then $(\forall n\in\N)\,P_n$.
\begin{proof} Let $F\defeq\{n\in\N:\neg P_n\}$ and suppose $F\neq\emptyset$. By the standing well-ordering assumption, choose the least $m\in F$.\\
By (i), $1\notin F$, so $m\neq1$. Hence $m-1\in\N$ and $m-1<m$. Minimality gives $m-1\notin F$, so $P_{m-1}$ is true.\\
Applying (ii) at $n=m-1$ gives $P_m$, contradicting $m\in F$. Thus $F=\emptyset$, so $(\forall n\in\N)\,P_n$.\\
\hlt{Structural Commentary}: Well-ordering supports a least-counterexample contradiction.\\
\hlt{Key Insight}: A least failure has a smaller successful predecessor.
\end{proof}
\end{theorem}
